% 第1章 1.2.2：连续斜坡 g_n 在 L^2([-1,1]) 误差下逼近阶跃 sgn。
% g_n(x)=-1 (x<=-1/n)，nx (|x|<1/n)，1 (x>=1/n)；画 n=2,4,16。
% 阴影为 n=2 时 |g_2-sgn| 所在区域；||g_n-sgn||_2^2 = 2/(3n)。
\begin{tikzpicture}[font=\figurefont\footnotesize]
\begin{axis}[width=82mm,height=46mm,axis lines=middle,
  x axis line style={draw=FigureStroke,line width=.6pt,-{Stealth[length=2mm]}},
  y axis line style={draw=FigureStroke,line width=.6pt,-{Stealth[length=2mm]}},
  tick style={draw=FigureStroke},tick label style={font=\figurefont\scriptsize,text=FigureInk},
  xmin=-1.1,xmax=1.18,ymin=-1.3,ymax=1.35,xtick={-1,-0.5,0.5,1},xticklabels={$-1$,$-\tfrac12$,$\tfrac12$,$1$},
  ytick=\empty,xlabel={$x$},xlabel style={anchor=north west,font=\figurefont\scriptsize},
  legend style={draw=none,fill=none,font=\figurefont\scriptsize,cells={anchor=west},at={(0.02,0.98)},anchor=north west},
  legend cell align=left]
\addplot[fill=FigureBlueFill,draw=none,forget plot] coordinates {(0,0) (0.5,1) (0,1)} -- cycle;
\addplot[fill=FigureBlueFill,draw=none,forget plot] coordinates {(0,0) (-0.5,-1) (0,-1)} -- cycle;
\addplot[draw=FigureMuted,line width=1.1pt,dashed] coordinates {(-1,-1) (0,-1)};
\addlegendentry{阶跃 $\operatorname{sgn}$}
\addplot[draw=FigureMuted,line width=1.1pt,dashed,forget plot] coordinates {(0,1) (1,1)};
\addplot[draw=FigureYellow,line width=1.1pt] coordinates {(-1,-1) (-0.5,-1) (0.5,1) (1,1)};
\addlegendentry{$g_2$}
\addplot[draw=FigureGreen,line width=1.1pt] coordinates {(-1,-1) (-0.25,-1) (0.25,1) (1,1)};
\addlegendentry{$g_4$}
\addplot[draw=FigureBlue,line width=1.1pt] coordinates {(-1,-1) (-0.0625,-1) (0.0625,1) (1,1)};
\addlegendentry{$g_{16}$}
\node[font=\figurefont\scriptsize,text=FigureInk,anchor=east] at (axis cs:-0.03,1) {$1$};
\node[font=\figurefont\scriptsize,text=FigureInk,anchor=west] at (axis cs:0.03,-1) {$-1$};
\end{axis}
\end{tikzpicture}
